\subsection{THETA3-C31: higher-CM kernels from horizontal homomorphism-lifting cycles on an actual boundary}\label{C31}

The elliptic-power restriction in C29 can be removed at a specified
good prime if one allows higher-codimension boundary cycles and a
polarization contraction. The cycles below are actual loci where
a specified homomorphism lifts to the tautological semiabelian
extension. Their realization as closures of interior special cycles
is a further comparison, not part of the assertion.

\paragraph{Data and assertion.}
Let $B/R$ be an abelian scheme of relative dimension $d$, where
$R$ is a characteristic-zero discrete valuation ring with residue
characteristic $p$. Suppose its generic fibre has rank-one CM by
the maximal order $\mathcal O_E$, and fix an $E$-compatible
polarization $\lambda$ of degree prime to $p$. Assume $p$ is
unramified in $E$. Fix a rigidified line bundle $Q$ inducing
$\lambda$. Give $B^\vee$ the conjugate-dual action
$a\mapsto\iota_B(\bar a)^\vee$, so that $\lambda$ is
$E$-linear; character-lattice actions below use the same
involution convention as Lan. These hypotheses hold for the chosen CM data outside
an explicitly retained finite set of primes. Neither a prime
dividing the polarization degree nor a ramified PEL prime is
silently included.

Let $P_\lambda$ be the normalized Poincar\'e kernel on $B\times B$
and, for $a\in\mathcal O_E$, put
\[
 P_a=([a]\operatorname{pr}_1,\operatorname{pr}_2)^*P_\lambda.
 \tag{C31a}
\]
On $B^2\times B$, with coordinates $(P,Q',T)$, define
\[
 \begin{aligned}
 j_+(P,Q')&=[a]P+Q',\\
 j_1(P,Q')&=[a]P,\qquad j_2(P,Q')=Q',\\
 \Gamma_j&=\{T=j(P,Q')\}.
 \end{aligned}
 \tag{C31b}
\]
Each $\Gamma_j$ is smooth over $R$, is a regular immersion of
codimension $d$, and is a labelled component of an actual
horizontal homomorphism-lifting locus on a prime-to-$p$ isogeny
pullback of a toroidal boundary stratum of a larger unitary PEL
model. If $\pi_{12}:B^3\to B^2$ is projection, then
\[
 c_1(P_a)=\pi_{12*}\left(
   ([\Gamma_{j_+}]-[\Gamma_{j_1}]-[\Gamma_{j_2}])
                  \cdot\operatorname{pr}_3^*c_1(Q)\right)
                  \quad\hbox{in }\CH^1(B^2).
 \tag{C31c}
\]
This is an integral Chow identity on the indicated pullback.
Its graph construction has no factor $24$. It does not provide
the separate rational theta section of C29. Descent from the
isogeny pullback retains its actual degree; that degree is a
unit at the specified prime.

\paragraph{The extension obstruction, as a scheme-theoretic locus.}
For an abelian scheme $A/S$, the normalized Poincar\'e bundle
on $A\times A^\vee$, regarded as a $\mathbb G_m$-torsor,
is the universal extension of $A$ by $\mathbb G_m$ over
$A^\vee$. The addition law in its first variable gives the
group law. The relative Picard universal property and the
rigidification identify its class with the parameter itself.
This is the relative form of Milne, \emph{Abelian Varieties},
Propositions 9.2--9.3, printed p.~41; its biextension and
base-change formulation is SGA~7~I, VII~2.9.5, p.~175,
and VIII~3.1--3.2, pp.~254--255.

More generally an extension $G$ of $A$ by a split torus with
character lattice $X$ is classified by
$c:X\to A^\vee$: push out by each character and apply the
preceding construction. For $h:B\to A$, a lift
$\widetilde h:B\to G$ exists exactly when the pulled-back
extension is split, equivalently
\[
                        h^\vee c=0.                 \tag{C31d}
\]
The assertion is functorial under every base change, so it
defines the actual zero fibre of the morphism in (C31d),
including its scheme structure. A splitting is unique because
$\Hom(B,\mathbb G_m)=0$, hence $\Hom(B,T)=0$ for a split
torus. One can verify the first equality directly: a map
from the proper geometrically connected abelian scheme to
the affine group factors through the base, and a group map
must preserve zero. If $G$ and $h$ carry compatible
$\mathcal O_E$-actions, this uniqueness also makes the lift
$\mathcal O_E$-linear: the two composites with any element
of the order lift the same homomorphism and therefore agree.

\paragraph{An actual unitary boundary carrying these extensions.}
Use the zero-dimensional unitary PEL datum of $A=B^3$ with
its product polarization and its full CM type. Take the
rank-one $\mathcal O_E$-lattice $\mathcal Q=\mathcal O_E$
in Lan, \emph{Toroidal compactifications of PEL-type Kuga
families}, corrected author version, Section~3A,
pp.~27--30. Its construction adds the two dual isotropic
lattices to $H_1(A)$ and gives a larger PEL datum with a
rational boundary. A top-dimensional cone gives an actual
boundary abelian stratum, prime-to-$p$ isogenous to
$\Hom_{\mathcal O_E}(\mathcal Q,A)=A$. The neat level
conditions in that section make it an abelian scheme,
rather than an unspecified torsor. The tautological
semiabelian extension on this stratum is described in
Section~3B, p.~31, and in the paragraph before Lemma~4.9
and its proof, p.~46. Its extension parameter is
$c:X\to A^\vee$, with $X$ of rank one over $\mathcal O_E$.
The order action and the polarization are the given ones.

This use includes the author's errata items (4), (8), (13),
and (23): the finite \'etale cover of the smaller
moduli space, the actual torus-torsor chart, both level
conditions and the reduction to principal level are retained.
No directedness assertion about the set of compactifications
is needed. The unramified and prime-to-$p$ polarization
hypotheses above place the construction in Lan's good-prime
base. After acquiring a sufficiently fine prime-to-$p$ level,
pull the boundary stratum to the specified CM base $R$.

Here is why the isogeny ambiguity does not change the claimed
labelled graphs. Write $C_\Phi$ for the boundary parameter
abelian scheme. Forgetting the auxiliary parameter and level
defines
\[
 F:C_\Phi\longrightarrow
       \Hom_{\mathcal O_E}(X,A^\vee)^\circ.
\]
This is an actual prime-to-$p$ isogeny. Indeed the polarization
relation is $c\phi_{\rm cusp}=\lambda_Ac^\vee$. The chosen
cusp lattice map and $\lambda_A$ are isogenies prime to $p$;
solving this equation for $c^\vee$, and acquiring the specified
prime-to-$p$ level lifts, are finite \'etale isogenies.
Taking the identity component preserves surjectivity onto the
connected parameter variety. The component and pairing
conditions are precisely those in Lan's construction and its
corrected principal-level description in Lemma~4.9.

A rank-one projective lattice $X$ is free locally at all
primes over $p$. Choose a generator at those finitely many
primes and clear denominators only away from $p$. This gives
an $\mathcal O_E$-linear injection $i:X\hookrightarrow
\mathcal O_E$ of index prime to $p$. Define explicitly
\[
        \varphi(P)(x)=i(x)\cdot\lambda(P).
\]
The dual action is the conjugate-dual one fixed above.
This formula and the prime-to-$p$ degree of $\lambda$ give
a prime-to-$p$ isogeny
\[
       \varphi:B\longrightarrow
               \Hom_{\mathcal O_E}(X,B^\vee)^\circ.
\]
For $A=B^3$ take its third power. Choose an integer $m$
prime to $p$ annihilating $\ker F$.
Then $[m]F^{-1}$ is an actual homomorphism. Composing with
$\varphi^3$ supplies an actual isogeny pullback on which
the extension parameter is $[m]\varphi^3(P,Q',T)$.
All finite indices and isogenies in this construction are
recorded choices, not integral inverse isogenies.

For any of the three maps $j$ in (C31b), let
$h_j:B\to B^3$ be the polarization adjoint of
\[
                 (P,Q',T)\longmapsto j(P,Q')-T.
\]
It is an actual $\mathcal O_E$-linear homomorphism.
For example the first is the column $(\bar a,1,-1)$;
the other two are $(\bar a,0,-1)$ and $(0,1,-1)$.
The product polarization and Rosati conjugation verify
these formulas directly. Equation (C31d) on the isogeny
pullback is now
\[
               [m]\varphi\bigl(j(P,Q')-T\bigr)=0.
 \tag{C31e}
\]
The kernel of $[m]\varphi$ is finite \'etale over
$R$. After a finite \'etale extension supplying its
labels, (C31e) is a disjoint union of translates of the
graph $T=j(P,Q')$. Its zero-labelled component is exactly
$\Gamma_j$. In particular it is flat and smooth, with
no vertical component or hidden intersection multiplicity.
This proves the boundary and homomorphism-lifting assertions.
The zero label is essential: the entire unlabelled
lifting locus is not silently replaced by this component.

\paragraph{The polarization contraction.}
The rigidified theorem of the square gives
\[
 P_a\simeq j_+^*Q\otimes(j_1^*Q)^{-1}
                         \otimes(j_2^*Q)^{-1}.
 \tag{C31f}
\]
If $i_j:B^2\hookrightarrow B^3$ is the graph immersion,
then $\pi_{12}i_j=\mathrm{id}$ and
$\operatorname{pr}_3i_j=j$. The projection formula for
the first Chern operator therefore gives
\[
 \pi_{12*}\bigl([\Gamma_j]\cdot
                  \operatorname{pr}_3^*c_1(Q)\bigr)
     =c_1(j^*Q)\cap[B^2].
 \tag{C31g}
\]
Taking the signed sum and using (C31f) proves (C31c).
Every graph is a section of the smooth morphism
$B^3\to B^2$, so it is a regular immersion of codimension
$d$ and is isomorphic to $B^2$. This proves its remaining
geometric properties without a transversality assumption.
The codimension count is $(d+1)-d=1$ after pushforward.

Apply (C31c) to the finite list of $a$'s in the CM
embedding projector of C23, retaining its algebraic
coefficients. It expresses every selected pure kernel of
A6 as a finite sum of these specified boundary
homomorphism-lifting cycles contracted with a polarization.
Substitution in C23's formula for $\mathcal C$ now works
for arbitrary rank-one CM factors satisfying the good-prime
hypotheses, not only elliptic powers. The first Chern
operations and proper pushforwards are exactly those in
that formula. For an arbitrary further pullback, use the
derived pullback of the regular graph's Koszul complex;
otherwise excess Tor terms could be lost. The Chow identity
itself requires no ordinary proper-intersection assertion.

\paragraph{Remaining comparison.}
These are actual labelled horizontal cycles on a specified
unitary boundary isogeny chart, with a precise moduli
meaning. They are of codimension $d$ before the
polarization contraction. When $d>1$, one cannot replace
their kernels by divisors by analogy with an elliptic
curve. The argument has not identified each boundary
lifting locus with the full formal completion of a
particular interior Kudla--Rapoport cycle. That requires
compatibility with the remaining degeneration data and
its generic-fibre special-homomorphism condition. Nor
has it supplied an interior correspondence from the
original compact 032 ball to the larger model. Those
are the remaining comparisons; the higher-CM boundary
host and its horizontal kernel contraction are now
specified.

\framework{The universal extension and its splitting
obstruction are Poincar\'e/biextension constructions in
SGA~7~I, VII~2.9.5, p.~175, and VIII~3.1--3.2,
pp.~254--255. Lan's Section~3B uses Raynaud's
semiabelian quotient and extension results in SGA~7~I,
IX~1.4 and IX~2.4. Formula (C31g) is an explicit
divisorial correspondence and pushforward of the kind
recorded in cote~44, scans~7--8; it uses C23's exact
first-Chern and specialization framework.}
